% !TeX program = xelatex
% Source-only research draft. No LaTeX/Lean/Python compilation or execution
% was performed for this document. Historical validation is explicitly scoped.
\documentclass[UTF8,a4paper,11pt,fontset=fandol]{ctexart}
\usepackage[margin=25mm,headheight=14pt]{geometry}
\usepackage{amsmath,amssymb,amsthm,mathtools}
\usepackage{booktabs,longtable,array,enumitem}
\usepackage{xcolor,listings,fancyhdr}
\usepackage[unicode,colorlinks=true,linkcolor=blue!55!black,
  citecolor=blue!55!black,urlcolor=blue!55!black]{hyperref}
\usepackage{microtype}
\setlength{\parindent}{2em}
\setlength{\parskip}{0.25em}
\setlist{nosep,leftmargin=2em}
\setlength{\emergencystretch}{3em}
\renewcommand{\arraystretch}{1.18}
\pagestyle{fancy}
\fancyhf{}
\fancyhead[L]{树上精确路由：成果边界与研究推进}
\fancyhead[R]{2026-11-18 草稿}
\fancyfoot[C]{\thepage}
\newtheorem{theorem}{定理}[section]
\newtheorem{proposition}[theorem]{命题}
\newtheorem{lemma}[theorem]{引理}
\newtheorem{corollary}[theorem]{推论}
\theoremstyle{definition}
\newtheorem{definition}[theorem]{定义}
\newtheorem{example}[theorem]{例}
\newtheorem{remark}[theorem]{说明}
\newcommand{\N}{\mathbb N}
\newcommand{\Z}{\mathbb Z}
\newcommand{\Q}{\mathbb Q}
\newcommand{\R}{\mathbb R}
\newcommand{\supp}{\operatorname{supp}}
\newcommand{\ic}{\operatorname{intcone}}
\newcommand{\parv}{\operatorname{pa}}
\newcommand{\ch}{\operatorname{ch}}
\newcommand{\cost}{\operatorname{cost}}
\newcommand{\bits}{\operatorname{bits}}
\newcommand{\Feas}{\operatorname{Feas}}
\newcommand{\Peak}{\operatorname{Peak}}
\newcommand{\Hist}{\operatorname{Hist}}
\newcommand{\canon}{\operatorname{canon}}
\newcommand{\repo}[1]{\nolinkurl{#1}}
\newcommand{\status}[1]{\textbf{证据层级：#1。}}
\lstset{basicstyle=\ttfamily\footnotesize,breaklines=true,
  columns=fullflexible,keepspaces=true,showstringspaces=false,
  frame=single,rulecolor=\color{black!20},backgroundcolor=\color{black!2},
  literate={·}{{\textbullet}}1}

\title{树上精确边流的可认证预算优化\\
\large 当前成果、稀疏证书、接口下界与置换商动态规划}
\author{研究讨论稿}
\date{2026 年 11 月 18 日}
\begin{document}
\maketitle

\begin{abstract}
本稿整理固定树、精确边访问重数、有符号逐次服务费用与个体预算下的路线分解研究，
并提出可以接续现有 Lean 开发的三个研究包。现有系统已经建立路线与边数矩阵的
精确对应，给出完整证书搜索、两条不同的执行计费链、最小预算与重复预算查询，
以及矩阵到显式顶点路线的重建。本文不把这些不同层次拼成一个未经证明的
统一机器性能结论。研究推进的重点是三个可独立检验的数学问题：一是把单路线
可行域直接线性化，并证明压缩模式证书同时满足对数支持界和奇偶支持界；
二是用五个顶点的非负费用实例证明接口必须保留足够的逐路线费用信息，
并明确此类下界只针对等价类、显式代表及指定兼容矩阵；
三是通过路线置换、整数运输表和完整见证提升，给出商动态规划的正确合并方法。
本文同时区分已有 Lean 定理、本文纸面证明、借助外部文献得到的推论及待实施算法，
避免把已知的整数锥、对称消除或秩压缩思想当作本项目的新发明。
源码按约二十页的研究草稿组织，尚未排版编译，实际页数取决于本地字体与工具链。
\end{abstract}

\noindent\textbf{阅读与版本约定。}
仓库为 \url{https://github.com/huam54925-cpu/uav-routing-certification}，
研究分支为 \repo{research/structural-service-cuts-20261023}。
本文核对的远端基线为
\repo{f3ea0823be69ffa7152dec903069c2ba79ad20a4}。
本文新增的是研究论证和计划，不包含新的 Lean 构建、Python 实验或公理审计结果。
文中出现的历史 PASS、测试数量和计费定理均保留其原记录的执行器、日期与假设。
所列 Lean 片段仅是未编译的接口设计，不以占位证明冒充形式化成果。

\tableofcontents

\section{研究问题与可以主张的贡献}

项目最初来自多无人机巡检算例，但目前最成熟的研究对象已经是一个独立的
组合优化问题：给定一棵根树和每条边必须实现的精确访问重数，能否把全部访问
分给指定数量的合法 depot 路线，使每条路线的服务费用不超过同一预算？
服务费用按每次访问累加，允许为负；路线可以重复访问内部顶点，但不能在中途
返回 depot 后再开始另一次出发。这些条件共同决定了数学问题，不能在文献比较中
只保留“树上车辆路由”这一名字。

现有项目值得强调的部分，是同一语义下多种表示之间的严密对应：路线、每路线
边数向量、矩阵、带重数的模式集合、搜索执行器、检查器和最终输出。
这种对应使后来更换算法时有一个稳定的语义参照。它既有研究价值，也为新的
结构定理提供基础：如果能减少必须保留的模式数、证明某种接口无法进一步压缩，
或者给出正确的对称商运算，就不只是增加验证代码，而是在解释问题本身。

另一方面，完整性不能只按模块数量衡量。带地址计费的旧全局搜索与标量计费的
子树 DP 是不同执行器；保存旧枚举次序所选证书与只求任一最优证书，是不同规格；
生成二进制压缩证书与展开所有路线顶点，是不同输出任务。本文把这些分界写入
定理与计划，而不把它们留给读者猜测。

本文建议形成如下研究主张。第一，固定精确边流的可行性具有一个低维整数锥表示，
并且总有编码长度受输入位长多项式控制的模式证书。第二，即使树只有五个顶点、
费用全部非负，某些自然接口的补全等价类也可能随二进制数值指数增长；
因此压缩必须说明保留什么语义。第三，车辆编号具有可利用的对称性，但正确
去掉编号后，子树合并必须显式处理相对匹配；整数运输表提供了可证明、可重建的
表达。这三点分别回答“证书可以多短”“接口哪些信息不能丢”“搜索如何去掉
冗余编号”，比单独增加一层内存计费更容易形成清晰的新增贡献。

\section{模型、访问语义与矩阵刻画}\label{sec:model}

\subsection{根树与精确边流}

设 $R=(V,E,o)$ 是有限根树，$o$ 为 depot。
对 $v\neq o$，记父顶点为 $\parv(v)$、孩子集合为 $\ch(v)$。
输入 $H:V\setminus\{o\}\to\N$ 指定边
$\{\parv(v),v\}$ 在两个方向上各被使用 $H_v$ 次。
因此 $H_v$ 不是一条边是否必须访问的布尔量，也不是优化变量；
解必须精确实现输入的每个数值。
令全部非根顶点数为 $n=|V\setminus\{o\}|$，并记
\[
M=\sum_{v\neq o}H_v,\qquad
A_H=\{v\neq o:H_v>0\},\qquad t=|A_H|,\qquad
U=\max(\{H_v:v\neq o\}\cup\{0\}).
\]
路线数为 $m\in\N$，服务费用 $s_v\in\Q$，统一预算为 $T\in\Q$。
本文按仓库的非 depot 费用约定令 $s_o=0$；若应用要求收取固定端点费用，
应先明确其计数规则并单独完成相应预算平移。

一条合法路线是有限顶点字
$r=(o,v_1,\ldots,v_{\ell-1},o)$，相邻顶点沿树边移动，
内部顶点均不等于 $o$。路线不要求简单，也不允许把多次 depot 出发合成
一条路线来规避个体预算。整个路线族必须有恰好 $m$ 条路线，
并按两方向的实际边经过次数实现 $H$。
当 $m=0$ 时，空路线族只可能实现零边流；若实现，它对任何预算都可行。
因此最小预算 API 必须区分“所有预算可行”和“存在有限最小预算”，
不能把空族的峰值擅自定义成零。

\subsection{每次访问计费及线性化}

对路线 $i$，令 $p_i(v)$ 为父边沿向下方向的经过次数。
树的割性质保证向上经过次数相同。记 $n_i(v)$ 为路线对非根顶点 $v$ 的访问次数，
则每次到达 $v$ 要么来自父亲，要么来自一个孩子，所以
\begin{equation}\label{eq:visits}
n_i(v)=p_i(v)+\sum_{u\in\ch(v)}p_i(u).
\end{equation}
这个公式计入重复访问。例如路径 $o,h,x,h,o$ 在 $h$ 发生两次访问，
并非只支付一次 $s_h$。由有限求和换序可得
\begin{align}
C_i
 &=\sum_{v\neq o}s_v n_i(v)\notag\\
 &=\sum_{v\neq o}\bigl(s_v+s_{\parv(v)}\bigr)p_i(v)
   =\sum_{v\neq o}w_v p_i(v),\qquad
w_v:=s_v+s_{\parv(v)}.                    \label{eq:linear-cost}
\end{align}
这个等式对有符号有理数成立，且费用只依赖边数向量，不依赖内部合法访问顺序。
它解释了为什么矩阵证书足以认证这里的预算；对于与访问顺序、时间窗或连续轨迹
相关的费用，不能直接沿用同一个结论。

\begin{definition}[结构化边数矩阵]\label{def:matrix}
一个 $m$ 行矩阵 $P=(p_i(v))$ 可行，当且仅当以下四类条件同时成立：
\begin{align}
&p_i(v)\in\N,\qquad \sum_{i=1}^m p_i(v)=H_v;          \label{eq:columns}\\
&\sum_{v\in\ch(o)}p_i(v)=1 \quad(1\le i\le m);       \label{eq:departure}\\
&p_i(v)>0,\ \parv(v)\neq o\ \Longrightarrow\
    p_i(\parv(v))>0;                                \label{eq:closure}\\
&\sum_{v\neq o}w_vp_i(v)\le T \quad(1\le i\le m).    \label{eq:budget}
\end{align}
\end{definition}

条件 \eqref{eq:columns} 保证精确边流，\eqref{eq:departure} 保证每行恰好一次
depot 出发，\eqref{eq:closure} 保证正支持连接到 depot，
\eqref{eq:budget} 保证逐路线预算。由 \eqref{eq:columns} 自动得到
$0\le p_i(v)\le H_v$。这些条件组成仓库已经形式化的矩阵刻画，
而不是本文新提出的路线定义。

从路线到矩阵，只需计数并使用树边割的平衡；反向则对每行的正支持构造路线。
一种构造是在第一次进入正孩子边时处理其后代，再把该边剩余重数实现为往返。
祖先闭合保证每个必须处理的后代都能从根到达，而根行和等于一保证中途不回库。
由于式 \eqref{eq:visits}，重建保留访问数和任意有符号线性服务费用。
已有路线桥接还保留它在具体接口上需要的树数据和度数假设；
引用现有 Lean 定理时应携带这些假设，不能从纸面描述中省掉。

\subsection{模式重数、根分解与任务边界}

将相同的矩阵行合并，得到互异模式 $p^1,\ldots,p^k$ 及正整数重数
$\lambda_1,\ldots,\lambda_k$：
\begin{equation}\label{eq:patterns}
\sum_{j=1}^k\lambda_jp^j=H,\qquad
\sum_{j=1}^k\lambda_j=m.
\end{equation}
每个模式仍须满足结构条件和预算。这样压缩不会把单条路线的预算平均化，
因为每个模式都独立通过行检查。
模式证书可以在不立即展开 $m$ 行的情况下表达重复路线，适合二进制输入。
它与显式路线族的等价已经有仓库基础，本文的新问题是保证 $k$ 有多小，
而不是重新发明 \repo{pattern_counts}。

根的孩子把路线族分成不相交的分支：由根行和等于一，每条路线只能进入一个
根孩子分支。根分支分解在仓库中早已完成，主要接口包括
\repo{matrix_iff_branch_matrices}、
\repo{pattern_counts_iff_local_counts} 与 \repo{zero_branch_iff}，
记录于 \repo{TREE_BUDGET_ROOT_FACTORIZATION_STATUS.md}。
本文后续的接口与置换研究主要针对非根内部合并；不会把已有根分解列成新的成果。

从 UAV 应用到本模型，已经进行了明确的限制与抽象。网络必须是树或已经有独立
证明的树支持结构；边流 $H$ 固定；费用是每次访问的加和；路线数固定。
一般图的循环选择、时间窗、任务先后、动力学航迹和电池充放电均未由这些定理覆盖。
线性方向旅行费用可以通过已有虚拟服务费编码复用检查器，但编码可能产生负费用，
所以不能据此继续使用仅在非负条件下成立的剪枝。应用层的扩展需要新的保真定理，
而不是仅增加解释文字。

\section{截至基线的成果及其证据范围}\label{sec:status}

\subsection{全局搜索与子树 DP 是两条不同执行链}

旧全局搜索以完整残余向量及剩余路线数 $(k,r)$ 为状态。
在固定输入上定义
\[
S(H)=\prod_{v\neq o}(H_v+1).
\]
已有结果给出 $S(H)\le 2^M$ 与 $S(H)\le(U+1)^t$，
并把候选生成、数组向量、缓存顺序、完整答案和显式输出连接到带地址存储的执行器。
\#22 的计费模型下已有
\begin{equation}\label{eq:old-cost}
\mathrm{steps}+\mathrm{allocations}
\le 2^{50}S(H)^3(\ell+1)^7
\end{equation}
及相应的 FPT$(M)$、FPT$(t,U)$ 结论，$\ell$ 表示其输入编码长度。
该结果有明确的地址、操作和分配语义，覆盖声明中的输入字符串处理，
但不是编译后 Lean 程序的墙钟时间，也不是未经另外构造就自动得到的标准
图灵机模拟定理。

\#19 的子树 DP 沿树合并带路线标签的局部状态。其接口保留逐路线活跃信息与费用，
并携带完整矩阵见证；LCM 缩放把有理费用与预算转换为等价的整数执行。
H 包在创建路线维向量前检查根总量，因此对如 $m=2^{200}$ 的不兼容输入，
拒绝代价依赖其二进制位宽，而不会先展开长度 $m$ 的向量。
I 包将子树 DP 的标量算术、坐标比较、列表控制及矩阵读加操作置于同一收费规则。
其中一项上界为
\[
C_{\mathrm{DP}}\le
4096(N+1)(m+1)^2(1+\mathrm{resourceWork}).
\]
H 和 I 使用同一保守的填充位宽；其中宽度包含操作上界本身，而非已经证明可用的
对数宽度。它完成了所选标量模型，并未声称紧输入位复杂度。

因此，新子树 DP 不能直接继承式 \eqref{eq:old-cost} 的完整地址成本结论。
它们可以共用语义和检查器，却有不同的递归、数据布局、共享策略与分配过程。
把“二者各自存在完整定理”拼成“新算法已经完成全机器成本精化”是不成立的。
原有工作块已经按各自范围交付；新存储层也不应被追加入旧包的验收条件，
从而使已完成任务永远处于未完成状态。

\subsection{最小预算、重复查询与断点}

固定 $R,H,s,m>0$ 后，可行矩阵只有有限多个。
对结构可行矩阵 $P$ 定义峰值
$B(P)=\max_i C_i$。若存在结构可行矩阵，则
\[
q=\min_P B(P),\qquad
\Feas(T)\ \Longleftrightarrow\ q\le T.
\]
仓库的 \repo{TreeBudgetSubtreeBudgetProfile.optimize}
已经给出最小有理预算及其证书，并处理空族的全预算可行情形。
它不是通过假设费用非负才能成立；有限候选集合与预算上闭性足以支持该结论。

最新查询接口 \repo{prepareQueries} 只准备一次根表，
随后 \repo{lookup T} 复用严格前缀极小值记录。
这个接口保留旧 \repo{solveAuto} 在预算 $T$ 下选中的完整证书，
包括费用并列与有序见证的选择，而不只是返回同一个布尔值。
查询计数统计测试了多少阈值，且不超过断点数；它不包含建表、缩放、
物化所选完整证书的全部成本。

在整数化权重下，令
\[
S=\sum_{v\neq o}w_vH_v,\qquad
W_+=\sum_{v\neq o}\max(w_v,0)H_v,
\]
并令 $g$ 为正质量坐标上权重的最大公因子。
已证断点 $b$ 满足 $S\le mb$、$b\le W_+$，
且在 $g>0$ 时 $g\mid b$。所以 $b/g$ 落在
\[
\left\lceil\frac{S}{mg}\right\rceil
\le \frac bg\le
\left\lfloor\frac{W_+}{g}\right\rfloor.
\]
当 $m=1$ 或 $g=0$ 时至多保留一个断点。
若规格要求保持原选中位置，每个被保留的位置都有必要。
这些是最新基线的已有成果，不应再包装成后续研究计划。

这里要区分两种问题：为了重现旧求解器在所有预算上的选择，需要保存一条有序
证书轮廓；为了仅返回某个最优解，只需要最小值及一份最优证书。
后续置换商算法优先采用第二种规格，才有空间真正合并标签不同的状态。

\subsection{堆去重、路线输出与验证边界}

\repo{TreeBudgetSubtreeHeapDedup.dedup} 已使用实际入口地址和关键字地址，
返回链共享原条目地址，并保留冲突中的最右代表。
若输入链长为 $n$、输出长为 $d$、标量轨迹中的完整后缀关键字映射数为 $K$，
则新增持久堆槽恰为
\[
3n+6K+3d.
\]
输入堆与共享 nil 哨兵在此前已存在，不能重复计入新增分配。
最新版 \repo{Work.dedupD} 进一步统计引用加载、节点创建与关键字读比较，
区分临时原语分配和持久节点槽，并证明良构输入上完整结果相同。
控制与观察计数器维护、其他 DP 阶段和全求解器地址时间仍是独立工作。

\repo{TreeBudgetSubtreeRouteOutput.vertices} 按证书顺序与重数输出顶点列表。
\repo{round_trip} 证明投影恢复规范化行并保留访问费用；
\repo{output_word} 连接检查通过的证书与合法、满足预算的路线字。
总输出恰有 $m$ 条列表、$2M$ 次边移动与 $2M+m$ 个存储顶点。
这既是输出规格，也指出二进制输入下展开可能指数长。
显式输出长度定理不是列表构造的运行时间定理；即使算法本身非常直接，
仍要指定实现和成本语义后才能主张后者。

仓库历史记录包括定向构建、有限回归、公理审计和独立枚举比对。
这些记录能支持对应版本和范围内的结论。未出现 \repo{sorryAx}
说明所审计声明没有依赖该占位公理，不能自动证明审计列表完整、
所有源文件冷启动可构建、编译后程序最快，或论文所有非形式化论证都正确。
小规模缓存实验中出现额外列表开销，也只能评价那个执行器和那些样例。
本文不把它上升为“所有认证算法都没有实用价值”，也不把有限回归上升为
通用性能保证。

\section{研究包 A：单模式线性化与两种稀疏证书界}\label{sec:sparse}

\subsection{直接有界多面体，无需引入支持布尔变量}

以下结果是本文给出的纸面证明，尚未在此次工作中加入 Lean。
把所有 $H_v=0$ 的坐标删去，其值固定为零；保留坐标数为 $t$。
为简化记号，仍把 $p$ 写成所有非根顶点上的函数，并约定被删坐标为零。
配套实施路线图先按全部非根顶点数 $n$ 写出保守界；
本节把固定零坐标显式消去，得到更细的 $2^t\le2^n$ 支持界。
这是对同一模型的维数精化，不改变旧参数 $t$ 的含义。
定义有界有理多面体
\begin{equation}\label{eq:polytope}
P_T=\left\{p\in\R^{A_H}:
\begin{array}{ll}
0\le p_v\le H_v & (v\neq o),\\
p_v\le H_v p_{\parv(v)}
 & (v\neq o,\ \parv(v)\neq o),\\
\sum_{v\in\ch(o)}p_v=1,&\\
w\cdot p\le T.&
\end{array}\right\}.
\end{equation}
这里含有已知整数系数 $H_v$，并无变量之间的乘积。
若一个正质量顶点的父边质量为零，相应约束将迫使该模式坐标为零；
随后总和等于 $H$ 的要求会检测不可行，不需要默默假定输入支持已经祖先闭合。

\begin{lemma}[整数点精确性]\label{lem:integerpoints}
$p\in P_T\cap\Z^t$ 当且仅当其零扩展是一个满足结构归一化和预算的合法矩阵行，
且逐坐标不超过 $H$。
\end{lemma}
\begin{proof}
只需证明乘数不等式与正支持祖先闭合在整数点上等价。
若 $p_v>0$ 而 $p_{\parv(v)}=0$，不等式给出
$p_v\le H_v\cdot0=0$，矛盾。
反过来，若祖先闭合成立：当 $p_v=0$ 时不等式由非负性成立；
当 $p_v>0$ 时，父坐标是正整数，因而至少为一，
所以 $p_v\le H_v\le H_vp_{\parv(v)}$。
其余约束逐项对应定义 \ref{def:matrix} 的单行条件。
\end{proof}

这个线性化利用了自然数间隙“正数至少为一”，故对分数点并不等价于
一个连续版本的支持逻辑。它的目的正是描述合法行的整数点集合，
并为整数锥算法提供一个凸、有界、二进制可描述的生成区域。
不需要把每个坐标拆成“是否活跃”和“流量”两个变量，
也不应在调用维数参数算法时平白把维数翻倍。

记 $\widehat P_T=P_T\times\{1\}$，
且对集合 $X$ 定义
$\ic(X)=\{\sum_{j=1}^k\lambda_jx_j:
k\in\N,\lambda_j\in\N,x_j\in X\}$，
空和为零。
\begin{theorem}[精确整数锥表示]\label{thm:intcone}
存在恰好 $m$ 条满足预算的精确树路线，当且仅当
\begin{equation}\label{eq:intcone}
(H,m)\in\ic(\widehat P_T\cap\Z^{t+1}),
\end{equation}
其中路线侧采用已有矩阵刻画所要求的树与输入兼容假设。
\end{theorem}
\begin{proof}
从矩阵出发，将相同的行归组得到式 \eqref{eq:patterns}；
引理 \ref{lem:integerpoints} 保证各个生成向量 $(p^j,1)$ 属于指定集合。
反向展开每个生成模式 $\lambda_j$ 次。
最后坐标保证总行数恰为 $m$，前 $t$ 坐标保证精确边流，
而每行由引理满足预算与结构条件。最后应用已有路线重建定理。
\end{proof}

最后坐标一在上述统一契约中显式固定模式个数，
使证书规格及外部整数锥定理的对接更直接。
它不是本问题不可消去的维数：在已有根总量守卫通过时，
对 $\sum_j\lambda_jp^j=H$ 的根孩子坐标求和，
再用每个模式的根行和为一，就得到
$\sum_j\lambda_j=\sum_{v\in\ch(o)}H_v=m$。
因此可以借助已有模式总数关系省去附加坐标；
本文保留它是统一表示上的便利，不主张必须增加一维。

\subsection{子集和碰撞给出的对数支持界}

\begin{theorem}[对数支持界]\label{thm:logsparse}
若实例可行，则存在互异模式的正重数证书，其模式数 $k$ 满足
\begin{equation}\label{eq:logsparse}
2^k\le(m+1)\prod_{v\in A_H}(H_v+1),\qquad
k\le\left\lfloor\log_2\left((m+1)
\prod_{v\in A_H}(H_v+1)\right)\right\rfloor.
\end{equation}
\end{theorem}
\begin{proof}
在所有可行模式证书中，选取支持数最小的一个，
记其互异增广生成向量为 $a_j=(p^j,1)$，正重数为 $\lambda_j$，
$j=1,\ldots,k$。考虑所有 $2^k$ 个子集 $I\subseteq[k]$ 的和
$a(I)=\sum_{j\in I}a_j$。由于 $\lambda_j\ge1$ 且生成向量逐坐标非负，
每个 $a(I)$ 的第 $v$ 坐标在 $0,\ldots,H_v$ 之间，
最后坐标在 $0,\ldots,m$ 之间。因此可取的整数向量至多为右侧盒子的大小。

若不同子集 $I,J$ 有相同和，先删除交集，得到不相交的
$I'=I\setminus J$ 和 $J'=J\setminus I$，仍有 $a(I')=a(J')$。
最后坐标说明 $|I'|=|J'|$，因此二者均非空。
令 $\delta=\min_{j\in I'}\lambda_j\ge1$，
把 $I'$ 上重数各减去 $\delta$，把 $J'$ 上重数各加上 $\delta$，
其余保持不变。新重数仍非负，至少一个原正重数变成零，
而因为两组生成向量之和相同，精确总向量 $(H,m)$ 不变。
没有引入新模式，每个保留下来的模式仍合法，故支持严格下降，
与最小性矛盾。因此所有子集和互异，得到所需计数不等式。
\end{proof}

该证明没有假设费用非负。非负的是生成向量 $(p,1)$，不是费用系数 $w$。
费用只用于判断生成向量是否属于 $P_T$，支持交换根本不改变单个模式。
也不需要把全部可行模式枚举出来才能证明存在短证书；
构造短证书的高效算法是另一个问题。

\subsection{凸性与奇偶中点给出的维数界}

\begin{theorem}[奇偶支持界]\label{thm:paritysparse}
若实例可行，则存在至多 $2^t$ 个互异模式的可行证书。
\end{theorem}
\begin{proof}
先使用显式 $m$ 行矩阵表示解。其每个坐标在有限区间
$0,\ldots,H_v$ 内，所以可行矩阵集合有限且非空。
从中选取使平方势
\[
\Phi(P)=\sum_{i=1}^m\|p_i\|_2^2
\]
最小的矩阵。若其中有超过 $2^t$ 个互异行，
按模二的 $t$ 维奇偶向量分类，必有不同的两行 $p,q$ 坐标逐项同奇偶。
它们的中点 $r=(p+q)/2$ 属于 $\Z^t$。
由于 $P_T$ 凸，$r\in P_T$；由引理 \ref{lem:integerpoints}，
$r$ 仍是合法预算行。

将这两个出现的行替换为 $r,r$，总行数与各列总量不变，
而
\[
\|p\|_2^2+\|q\|_2^2-2\left\|\frac{p+q}{2}\right\|_2^2
=\frac12\|p-q\|_2^2>0.
\]
这与平方势最小性矛盾。因而该矩阵互异行数不超过 $2^t$，
将其合并即可得到所需模式证书。
\end{proof}

直接在“支持最小”的证书上做中点替换，未必立即减少支持：
若两种模式原来都出现多次，替换各一次还可能引入新中点模式。
因此证明必须先用平方势得到一个具有奇偶界的可行矩阵，
不能把这个中间步骤省略。下面再用全局最小支持，才可将两个界同时写入
同一份存在性结论。

\begin{corollary}[双支持界]\label{cor:both}
每个可行实例都有模式证书，其支持数满足
\[
k\le
\min\left\{
2^t,\
\left\lfloor\log_2\left((m+1)
\prod_{v\in A_H}(H_v+1)\right)\right\rfloor
\right\}.
\]
\end{corollary}
\begin{proof}
取全局最小支持数 $k_{\min}$。
定理 \ref{thm:logsparse} 的证明直接约束 $k_{\min}$。
定理 \ref{thm:paritysparse} 提供一份支持至多 $2^t$ 的证书，
所以 $k_{\min}\le2^t$。取实现 $k_{\min}$ 的证书即同时满足两界。
\end{proof}

\subsection{编码长度与 NP 证明尚需补齐的部分}

令 $L$ 是明确二进制编码下的输入位长，包含树、$H,m$、
有理费用和预算的分子分母。记
\[
B_H=\bits(m)+\sum_{v\in A_H}\bits(H_v).
\]
在通常自定界编码的常数因子内，$B_H=O(L)$。
对数支持界给出 $k=O(L)$；一个模式逐坐标编码所需位数为
$O(\sum_v\bits(H_v))$，一个重数不超过 $m$，
因此模式与重数列表总长为 $O(L^2)$，再加上固定的分隔及维数信息仍为多项式。
这与显式矩阵可能有 $m$ 行、显式路线可能有 $2M+m$ 个顶点形成直接区别。

这项结果补上一般二进制实例存在短证书的核心数学环节，
但不能把它写成“本轮已经完成标准机器 NP 成员性形式化”。
还需指定证书的唯一解析、非法输入处理、有限自然数乘加、
有理不等式检查与中间位宽，并在选定机器或显式位操作模型中证明检查时间多项式。
检查器可以逐模式验证式 \eqref{eq:polytope}，
再累加式 \eqref{eq:patterns}；它不应为验证而展开 $\lambda_j$ 次。
费用检查可使用公分母整数化，但需要证明分母组合及中间整数位宽受 $L$ 的
多项式控制，而不是只写一个无成本的有理数比较。

已有原检查器与 \#22 的编码经验是可复用的组件，
却不自动验证这个新的压缩二进制格式。
合理的交付次序是先形式化两个有限组合证明及模式语义，
再复用已有检查器并增加二进制编码与位成本精化，
最后封装编码长度和计费定理。这样每个阶段都有独立、可审查的结果，
不会为了提前宣称复杂度类别而把尚未建立的机器桥接隐藏在术语中。

\section{研究包 B：五顶点非负实例的接口信息下界}\label{sec:interface}

\subsection{问题不是能否存一个总费用}

子树 DP 把一个局部解交给外部时，必须让外部判断两者能否拼接。
若摘要只保留总流量、总费用和正支持，而丢掉费用如何分配到各条路线，
外部的逐路线预算通常无法恢复。仓库已有具体分离反例；
本节进一步把这种现象组织成一个可任意放大的实例族，
从而计算补全等价类和兼容矩阵的规模。

本节区分三种对象。第一是带路线标签的局部解，其外部环境也按同一标签提供费用。
第二是允许重新排列路线标签的局部解，此时只要求存在某种合法配对。
第三是任意符号化编码，它可能用一个整数或公式紧凑地表示许多局部解。
前两者的代表集合下界，不等于第三者必须占用同样数量的存储位。
这个范围限制是下界本身的一部分。

\subsection{统一的树、流量和非负费用}

固定整数 $m\ge2$、$K\ge m$。令
\[
\mathcal C_{m,K}
=\{c=(c_1,\ldots,c_m)\in\Z_{>0}^m:\ \sum_i c_i=K\}.
\]
考虑五个顶点 $o,h,v,x,y$，根为 $o$，边为
$oh,hv,vx,hy$，其中 $h$ 的孩子是 $v,y$，$v$ 的孩子是 $x$。
规定
\begin{equation}\label{eq:five-flow}
H_h=H_v=m,\qquad H_x=K,\qquad H_y=(m-1)K,
\end{equation}
并取
\[
s_o=s_h=s_v=0,\qquad s_x=s_y=1,\qquad T=K.
\]
所有费用非负；变换后的边系数只有 $w_x=w_y=1$，其余为零。
沿 $hv$ 切开树，把 $v,x$ 视为内部，把 $o,h,y$ 视为外部，
并在接口记录每条路线的 $hv$ 边数。

对任意 $c'\in\mathcal C_{m,K}$，定义内部局部矩阵
$I(c')$：第 $i$ 行在 $v$ 上的边数为一，在 $x$ 上的边数为 $c_i'$。
对任意 $c\in\mathcal C_{m,K}$，定义外部环境 $E(c)$：
第 $i$ 行在 $h$ 上的边数为一，在 $y$ 上的边数为 $K-c_i$，
且要求接口在 $v$ 上的边数为一。
由于 $c_i\le K-m+1$，有 $K-c_i\ge m-1\ge1$；
故所有这些局部矩阵与环境都有相同的逐路线正支持和边界重数。

拼接后第 $i$ 行为
\[
p_i(h)=p_i(v)=1,\qquad p_i(x)=c_i',\qquad
p_i(y)=K-c_i.
\]
所有列和恰为式 \eqref{eq:five-flow}，根归一化与祖先闭合均成立。
内部费用为 $c_i'$，外部费用为 $K-c_i$，
所以完整行的费用为
\begin{equation}\label{eq:five-cost}
C_i=K+c_i'-c_i.
\end{equation}

\begin{theorem}[带标签接口的单位兼容矩阵]\label{thm:identity}
在上述同一实例族中，$I(c')$ 与 $E(c)$ 按标签拼接后满足预算，
当且仅当 $c'=c$。
因此以 $\mathcal C_{m,K}$ 同时索引内部状态和外部环境的兼容矩阵是单位矩阵，
阶数为 $\binom{K-1}{m-1}$。
\end{theorem}
\begin{proof}
由式 \eqref{eq:five-cost}，预算成立当且仅当每个 $i$ 都有 $c_i'\le c_i$。
而两个向量的坐标和同为 $K$，若至少一个坐标严格小于，
总和也严格小于，矛盾。因此所有坐标相等。
反向在相等时各条路线费用都恰为 $K$，自然可行。
正整数有序组成的个数由隔板计数为 $\binom{K-1}{m-1}$。
\end{proof}

定理并没有使用负费用，也没有隐藏一个随 $K$ 增长的图。
变化只发生在二进制编码的流量与局部行数值中。
因此它适合用于解释：即使边界拓扑极小，逐路线资源信息仍可能很丰富。
这与“树分解宽度小，所以资源接口一定很小”的直觉不同。

\subsection{允许路线置换后仍存在的分离}

令两个内部向量等价，当且仅当它们只差一个坐标排列。
允许内部与外部通过任意排列 $\pi\in S_m$ 配对。
预算条件变成
\[
c'_{\pi(i)}\le c_i\qquad(1\le i\le m).
\]
同样因为两侧总和相等，所有不等式必须取等。

\begin{corollary}[无标签接口下界]\label{cor:unlabelled}
允许任意路线配对时，$I(c')$ 与 $E(c)$ 可兼容，
当且仅当 $c'$ 和 $c$ 的坐标多重集相同。
取非递减正整数向量作为代表，得到阶数
$p_m(K)$ 的单位兼容矩阵，其中 $p_m(K)$ 是 $K$ 分成恰好 $m$ 个正部分的分拆数。
特别地，$m=2$ 时 $p_2(K)=\lfloor K/2\rfloor$。
\end{corollary}
\begin{proof}
前述逐坐标不等式与总和相等推出匹配后逐项相等，
这恰好是多重集相同；反向按相同值匹配即可。
当 $m=2$，非递减向量唯一写为 $(a,K-a)$，
$1\le a\le\lfloor K/2\rfloor$，故得到最后公式。
\end{proof}

固定 $m=2$ 时，实例位长 $\ell=\Theta(\log K)$，
所以该兼容等价类数和显式全代表集合的必要大小可以达到 $2^{\Theta(\ell)}$。
带标签版本的计数为 $K-1$，去掉标签后仍有 $\lfloor K/2\rfloor$，
说明路线对称性确能删除冗余，但不会使所有资源状态消失。
研究包 C 因此应当报告在哪些重复类型分布下压缩有效，
而不能承诺对所有实例都得到多项式大小的状态表。

\subsection{下界究竟排除了什么}

考虑一种对局部状态生成摘要的确定性方法，并要求相同摘要在所有允许的外部
环境下给出相同兼容答案。定理 \ref{thm:identity} 的任意两个不同内部向量，
都由其中一个专属环境分开。因此该方法至少需要相应数量的不同摘要值。
若进一步限定压缩只保留原状态的一个代表子集，且必须对所有外部环境保留
“存在可补全状态”这一语义，那么每个状态的专属环境迫使它自身被保留，
所以代表子集无法缩小。

同一个单位矩阵在任意域上秩均等于其阶数。
这阻止了对这个明确补全关系直接套用“连接矩阵自动低秩，所以总有小基”
的推论。现有 rank-based 方法依赖其目标连接关系的特殊代数结构，
不能只因为本项目也使用 DP，就声称存在同样的单指数代表族。

上述结论不排除一个参数化公式、生成函数或整数锥约束同时编码很多状态。
区分 $N$ 个等价类只必需 $\lceil\log_2 N\rceil$ 位标签，
不是 $N$ 位；本例本身用二进制整数 $a$ 就能描述
$(a,K-a)$。因此本文不主张“任意状态编码都要指数位空间”，
更不主张“所有求解算法都有指数时间下界”。
也不把这个具体矩阵的秩，偷换成任意重构后的连接关系都必然同秩。
研究价值在于准确排除一类过强的压缩希望，并为下一种压缩方式选择正确规格。

\subsection{正向拼接契约与待形式化接口}

只有反例不足以指导实现，还应当给出充分信息的拼接契约。
对当前矩阵 DP，一个自然接口至少同时携带逐路线边界重数或活跃标记、
逐路线累积费用，以及可把子表行映射到父表行的身份或匹配见证。
在单割边的树上，祖先闭合决定允许进入哪些父行，费用按同一匹配逐项相加；
内部其他结构由局部证书负责。这个接口是足够性的候选，
需要逐项对照已有 DP 的真实字段，而不能直接把一个过大的全矩阵称作
“最小接口”。

在 Lean 实施时，宜先定义“对所有给定外部见证兼容结果相同”的关系，
再证明五顶点族的分离与显式代表必要性。
若想证明更强的接口最小性，必须固定环境语言：外部能否使用负费用、
能否重排标签、能否增加路线、是否固定全局 $H$，都会改变等价关系。
本稿的族已经在非负费用、固定全局边流、固定路线数下工作，
所以不需要通过放宽这些条件获得分离。

\section{研究包 C：路线置换商与运输表合并}\label{sec:orbit}

\subsection{群作用与语义规格}

所有车辆遵循相同结构和同一预算，因此同时重编号所有矩阵行不改变可行性。
对 $\sigma\in S_m$ 定义
$(\sigma\cdot P)_i=P_{\sigma^{-1}(i)}$。
列和、根归一化、祖先闭合都在该作用下保持；
当 $m>0$ 时最大费用也保持，$m=0$ 则仍单独作为空矩阵分支。
这只是对称性起点，不足以证明可以在每个子树中独立排序，然后直接把同序行相加。
不同子树选择的规范编号通常没有共同语义，必须补上相对匹配。

设一个局部行的对外信息称为“联合类型”。
联合类型包含边界结构字段及费用，例如父边重数、必要的活跃信息和累积费用；
确切字段应由现有递推语义确定。
若某两个字段相关，必须由同一个行类型绑定。
分别排序活跃向量和费用向量会制造原来不存在的配对，
即使分别看两个多重集都没有丢失任何值，也可能破坏完整行的语义。

给定 $m$ 个局部行类型 $x_1,\ldots,x_m$，可用直方图
$\mu_a=|\{i:x_i=a\}|$ 表示其置换轨道。
但直方图只描述对外类型；内部矩阵仍须保留可恢复的见证或共享引用。
规范化的正确性至少应包含：每个原状态映射到同一轨道代表、
每个保留代表有具体合法见证、父状态生成的所有轨道与原带标签 DP 一致。
单独证明“排序幂等”不能替代这些义务。

\subsection{两个不能采用的压缩规则}

\begin{example}[独立排序后同序相加不完整]\label{ex:sort}
令两侧费用均为 $(0,10)$，且没有结构禁配。
同序配对得到 $(0,20)$，峰值为 $20$；交叉配对得到 $(10,10)$，
峰值为 $10$。因此只保留规范排序后的一次同序相加，会丢掉更好的完整解。
\end{example}

\begin{example}[当前最小峰值不是中间状态支配关系]\label{ex:peak}
同一边界上有两种局部费用状态
$A=(4,4)$ 和 $B=(0,8)$，其总费用都为八。
目前 $\Peak(A)=4<8=\Peak(B)$。
但若后续增加 $C=(0,8)$，并允许最佳匹配，则
\[
\min_{\pi}\max_i(A_i+C_{\pi(i)})=12,\qquad
\min_{\pi}\max_i(B_i+C_{\pi(i)})=8.
\]
所以删除当前峰值较大的 $B$ 会删除更好的未来解。
\end{example}

这两个反例都只有非负费用。它们说明失败来自逐路线费用的关联和补全方式，
并非只来自负费用。根表的严格前缀极小值是在全部费用已经累积完毕后使用，
不能因此直接在每个子树中只留最小峰值。
同样，在固定已处理精确列和的阶段，所有状态的费用总和常是固定的；
若两个费用向量逐坐标有支配且总和相同，就只能逐坐标相等。
因此常见 Pareto 规则在这里可能只剩重复删除，是否有效必须先分析守恒量。

\subsection{运输表精确刻画所有相对排列}

令 $A,B$ 分别为左右实际出现的有限联合类型集合。
给定 $x\in A^m$、$y\in B^m$，记重数为
$\mu_a=|\{i:x_i=a\}|$、$\nu_b=|\{j:y_j=b\}|$。
兼容谓词 $K(a,b)$ 只依赖联合类型，
合并类型函数为 $g(a,b)$；例如结构字段按父子规则更新，费用相加。
所谓“只依赖类型”是一条需要证明的充分性条件，
不能把仍依赖隐藏矩阵内容的检查器强行当成这样的函数。

\begin{definition}[可行运输表]
非负整数表 $N=(N_{ab})_{a\in A,b\in B}$ 可行，当
\begin{equation}\label{eq:transport}
\sum_bN_{ab}=\mu_a,\qquad
\sum_aN_{ab}=\nu_b,\qquad
\neg K(a,b)\Longrightarrow N_{ab}=0.
\end{equation}
它表示类型 $a$ 的多少行与类型 $b$ 的多少行相配。
\end{definition}

\begin{theorem}[相对排列与运输表的完整对应]\label{thm:transport}
可由合法相对排列得到的合并类型多重集，恰好是可行运输表得到的多重集：
\begin{align}
&\left\{\biguplus_{i=1}^{m}\{g(x_i,y_{\pi(i)})\}:
 \pi\in S_m,\ \forall i\ K(x_i,y_{\pi(i)})\right\}\notag\\
&\qquad =
\left\{\biguplus_{a,b}N_{ab}\cdot\{g(a,b)\}:
 N\text{ 满足式 }\eqref{eq:transport}\right\}.       \label{eq:orbit-convolution}
\end{align}
而且每张可行运输表都能提升为一个具体排列，进而提升为完整矩阵见证。
\end{theorem}
\begin{proof}
给定排列 $\pi$，令
$N_{ab}=|\{i:x_i=a,\ y_{\pi(i)}=b\}|$。
固定 $a$ 后，这些集合按 $b$ 分割所有类型为 $a$ 的左索引，
故行和为 $\mu_a$。由于 $\pi$ 是双射，固定 $b$ 的列和为 $\nu_b$。
禁配类型不会出现，所以得到式 \eqref{eq:transport}，
且按每对类型统计合并结果，得到等号左侧包含于右侧。

反向给定可行 $N$。把每个左类型索引集
$L_a=\{i:x_i=a\}$ 分割成互不相交的块 $L_{ab}$，
使 $|L_{ab}|=N_{ab}$；这是由行和条件保证的有限集合分割。
把每个右类型索引集 $R_b=\{j:y_j=b\}$
分割成块 $R_{ab}$，大小仍为 $N_{ab}$，由列和条件保证。
对每对 $(a,b)$，选择一个双射
$\phi_{ab}:L_{ab}\to R_{ab}$。
所有左块分割 $[m]$，所有右块也分割 $[m]$，
所以这些双射的并组成全局排列 $\pi$。
若 $i\in L_{ab}$，则 $x_i=a$、$y_{\pi(i)}=b$；
出现的块非空时 $N_{ab}>0$，禁配条件保证 $K(a,b)$。
由此实现所给运输表，并得到等号反向包含。

最后对右子问题的完整矩阵、payload 引用和需要重编号的字段应用同一排列，
再使用原带标签合并规则。由于局部合法性在统一行置换下保持，
而兼容性只依赖已验证的联合类型，所得矩阵是该合并直方图的具体合法见证。
这一步要求重排整行证书，不能只重排费用向量。
\end{proof}

这个组合事实可解释为稳定子双陪集与列联表的经典对应。
本项目可以贡献的是：确定当前 DP 的充分联合类型，
实例化结构兼容性，证明与原执行器的轨道级完整等价，
并把运输表实现提升成原检查器可接受的证书。
仅把上述经典对应重新命名，不足以构成算法原创性。

\subsection{状态数、匹配数及参数含义}

\begin{proposition}[直方图和运输表的计数]\label{prop:counts}
若一个阶段可能出现 $d\ge1$ 种联合行类型，则长度 $m$ 的直方图数为
$\binom{m+d-1}{d-1}$。
若某次左右状态分别实际出现 $a,b\ge1$ 种类型，
则给定边际的运输表数量至多
$(m+1)^{(a-1)(b-1)}$。
\end{proposition}
\begin{proof}
第一式是非负整数方程 $n_1+\cdots+n_d=m$ 的隔板计数。
第二式中，给定左上 $(a-1)\times(b-1)$ 个条目后，
前 $a-1$ 行的最后一列由行和唯一确定；
前 $b-1$ 列的最后一行由列和唯一确定；
右下角再由最后一行的行和确定，最后一列之和自动一致。
所以限制到左上块是单射。每个条目在 $0,\ldots,m$ 之间，
故最多有所述数量的块。某些块产生负数或违反禁配，只会减少数量。
\end{proof}

带标签类型向量的粗界为 $d^m$，规范直方图能够删去大量重编号。
具体轨道大小为 $m!/\prod_an_a!$；
当所有行同型时轨道只有一个元素，不能宣称普遍节省 $m!$ 倍。
左右各两类时，运输表只由一个整数决定，至多 $m+1$ 种，
比直接枚举所有相对排列更合适。
但如果每一行都有不同费用，类型数仍可能接近 $m$，
运输表枚举的最坏界并不小。

固定 $d$ 时直方图数关于数值 $m$ 为多项式。
只有在 $m$ 显式展开或按一元计入输入规模时，
形如 $m^{O(d^2)}$ 的界才给出相应 XP 表述；
即使如此也不能将其说成 FPT$(d)$。
当 $m$ 为二进制输入时，这只是关于数值参数的界，
不能据此得出关于二进制输入长度的 XP，
因为“关于 $m$ 的多项式”仍可能对输入位长是指数的。
若输出要求展开 $m$ 行，这一数值依赖有输出长度背景；
若输出是压缩模式证书，则应继续争取避免直接依赖数值 $m$。

对整数化后的子树费用，令
\[
W_v^+=\sum_{u\text{ 属于已处理区域}}H_u\max(w_u,0),\qquad
W_v^-=\sum_{u\text{ 属于已处理区域}}H_u\max(-w_u,0).
\]
每行局部费用位于 $[-W_v^-,W_v^+]$。
若边界结构字段只有 $b_v$ 种可能，则费用整数网格上的粗类型数为
$d_v\le b_v(W_v^-+W_v^++1)$。
这些界使用绝对正负质量，不能用可能严重相消的总费用替代。
它们可用于第一版可证明上界，但论文还应报告实际出现类型数、
中间直方图数及匹配数，判断理论压缩是否转化成实际改进。

\subsection{终层合并的更小算法与其严格适用范围}

若一次合并后直接检查最终预算，没有未来费用要再加到各行上，
给定左右直方图时可以只判断是否存在满足
$K(a,b)$ 且 $c(a)+c(b)\le T$ 的运输表。
构造源到左类型容量 $\mu_a$、右类型到汇容量 $\nu_b$ 的网络，
允许类型对之间加容量至少 $m$ 的边。
整数最大流值为 $m$ 当且仅当所需运输表存在，
由定理 \ref{thm:transport} 可恢复匹配证书。
若不可行，对应割提供该次固定直方图合并的不可行解释；
它并不自动证明全局所有其他直方图都不可行。

在完全没有结构禁配的纯费用情形，还有如下经典重排结论。
\begin{proposition}[终层反序匹配]\label{prop:antitone}
设 $m>0$。对任意实数向量 $x,y\in\R^m$，包括有符号费用，
\[
\min_{\pi\in S_m}\max_i(x_i+y_{\pi(i)})
=\max_i\bigl(x_{(i)}+y_{(m+1-i)}\bigr),
\]
其中括号表示非递减排序。
\end{proposition}
\begin{proof}
固定 $x$ 非递减。若 $i<j$ 对应的两个 $y$ 值满足 $u\le v$，
原两和为 $x_i+u,x_j+v$，最大者为 $x_j+v$。
交换以后 $x_i+v\le x_j+v$、$x_j+u\le x_j+v$，
所以最大值不会增加。反复交换可把 $y$ 排成非递增，得到一个最优排列。
证明只用了实数序与加法单调性，无需非负性。
\end{proof}

这个命题不是中间层的贪心正确性；例 \ref{ex:peak} 已说明，
为了未来补全，可能必须保留当前不是最优的费用分布。
它可作为一个有明确终止条件的局部加速器，或用于特殊 hub 子类，
其调用点必须在实现中显式证明“没有未来资源追加”。

\subsection{规范见证与旧选择顺序}

新算法建议提供 \repo{optimizeOrbit}：
返回与旧求解器相同的最小预算和一份通过原检查器的完整证书，
并记录可重复的规范匹配策略。
它不自动满足“每个预算下与旧求解器选中完全同一份证书”的更强规格。
旧执行器的最右重复代表和首个成功位置都依赖原枚举次序，
两个同轨道见证可能在这个次序中有不同的先后。

若未来确实需要原样重现旧选择，还须证明一个可组合的 provenance 顺序：
商状态保存的信息要足以恢复所有被折叠候选中最早满足旧规则的来源。
这个额外信息可能抵消一部分状态压缩。
因此验收应先区分“最优值相等”“证书有效”“规范输出稳定”
和“旧证书逐字段相等”四个不同目标。
保持旧执行器作为参照是合理的；要求任何创新算法都立即保持它的全部偶然顺序，
反而会把算法空间限制在实现细节上。

\section{文献定位：哪些是继承，哪些值得继续证明}\label{sec:literature}

\subsection{整数锥稀疏性不是本项目独有的方法}

Eisenbrand 与 Shmonin 的整数锥 Carath\'eodory 型结果
\cite{eisenbrand2006} 是稀疏生成表示的重要背景。
本文的子集和碰撞与奇偶中点证明属于这类经典思路的直接应用，
不主张首次提出这些一般方法。本文可检验的增量是：
精确构造适合当前路线语义的 $P_T$，
说明自然数祖先闭合无需额外布尔维数，
给出依赖实际总目标 $(H,m)$ 的短证书界，
并计划把这些结论连接到已有检查器与二进制压缩证书。
如果这些连接最后只是几条已知结论的机械转述，
论文中也应相应降低原创性主张；形式化与模型适配的价值应分别解释。

Goemans 与 Rothvo\ss{} 的工作
\cite{goemans2014,goemans2020} 研究常数物品类型的高重数装箱，
其 SODA 2014 版本 Theorem 2.2 给出的整数锥与多面体相交算法，
时间可写为
\[
\operatorname{enc}(P)^{\,2^{O(d)}}\,
\operatorname{enc}(Q)^{O(1)},
\]
并有仅依赖维数的支持控制。
将维数固定即可得到多项式时间，但输入位长的指数依赖 $d$，
因此这个式子属于 XP 形式，不能仅因为“对每个固定维数多项式”
就称为 FPT$(d)$。它非常适合给本项目提供第一层理论解释，
却不等于旧 \#22 的 FPT$(t,U)$ 已经自动改善为 FPT$(t)$。

\subsection{2026 年整数锥算法给出的新文献推论}

本次检索还核对了 Jansen、Ohnesorge 与 Schumacher
的预印本第 5 版（2026 年 11 月 3 日）
\cite{jansen2026convex}。该版本题名为
\emph{A Constructive FPT Algorithm for Integer Cone Optimization with Convex Costs}。
其 Theorem 1 在类型数 $\tau$、每类整数维数 $d$ 及文中定义的输入模型下，
给出形如
\[
2^{\,\tau\log(\tau+1)\,2^{O(d)}}\,|I|^{O(1)}
\]
的构造性算法；连续坐标维数在该定理的输入依赖中另有规定。
版本间标题和表述可能变化，因此引用应固定版本，而不是只引用搜索结果摘要。

对本文的固定预算判定，取一种生成类型 $\tau=1$、
没有连续坐标、凸成本函数为恒零，并使用增广生成域
$\widehat P_T$ 与目标点 $(H,m)$。
其整数维数为 $d=t+1$，生成域有界、所有约束有理，
矩阵系数及目标均由输入以多项式位长给出。
这些条件下得到以下\emph{依赖外部定理的纸面推论}：
\begin{corollary}[文献算法的路线特例]\label{cor:literature-fpt}
在所引用整数锥定理的输入与输出模型下，
固定精确树边流的预算判定可在
$2^{2^{O(t)}}L^{O(1)}$ 时间内完成，
并以压缩模式重数形式构造证书。
\end{corollary}
\begin{proof}
由引理 \ref{lem:integerpoints} 与定理 \ref{thm:intcone}，
路线判定等价于该整数锥的目标成员性。
生成域及目标的编码只增加输入位长的多项式量，
代入 $\tau=1,d=t+1$ 即得。
外部算法返回的生成向量均为合法模式，重数实现精确目标，
所以通过式 \eqref{eq:patterns} 取得压缩证书。
\end{proof}

这是文献定理对本模型的推论，不是当前仓库新实现的 FPT 算法，
也不是此次工作完成的 Lean 时间定理。
它与已有 \#22 结论并列时，应明确新旧算法不同、参数不同、输出形式不同。
若要求将证书展开为 $2M+m$ 个顶点，还必须加上不可避免的输出代价；
不能把压缩输出的输入位复杂度原样写到显式顶点输出上。
若把该方法进一步用于最小预算，需要另外给出有限有理搜索范围、
精度或断点选择及相应调用次数，本文不通过一句“二分即可”跳过这些义务。

已有根分支分解还能帮助分析子实例参数：
每条路线只属于一个根分支，可以在已证的分支接口下组织各分支判定。
但由于路线数在精确根边流下已经受约束，如何组合分支预算与压缩输出应直接
复用现有定理。这里的新工作是将文献算法按实际分支维数实例化并写清编码代价，
而不是再次声称发现了根分解。

\subsection{高重数装箱的困难性与双指数依赖}

高重数装箱给出一个很直接的模型嵌入。
设有 $d$ 种正整数物品大小 $a_j$、重数 $n_j$、箱容量 $B$，
要求使用至多 $m>0$ 个箱子。构造 depot $o$、hub $h$ 和叶子 $x_j$，
边为 $oh,hx_j$，取
\[
H_h=m,\quad H_{x_j}=n_j,\quad
s_h=0,\quad s_{x_j}=a_j,\quad T=B.
\]
每行由根归一化强制 $p_i(h)=1$，
叶子坐标 $p_i(x_j)$ 就是第 $i$ 箱中该物品的件数，
预算是 $\sum_ja_jp_i(x_j)\le B$。
空箱由 $o,h,o$ 表示，故“至多 $m$ 箱”与“恰好 $m$ 条允许无叶任务的路线”
一致。反方向从矩阵读出装箱即可。
忽略零重数类型后正边坐标数为 $t=d+1$，
映射不展开任何物品副本，只复制二进制重数及大小。

Jansen、Ohnesorge 与 Pirotton 的 ICALP 2026 论文
\cite{jansen2026lower} 给出高重数装箱的 ETH 条件下界，
其 Theorem 3 排除相应的
$|I|^{2^{o(d)}}$ 时间界。它也排除具有同样参数意义的
$2^{2^{o(d)}}|I|^{O(1)}$ 形态，因为该形态包含于被排除的范围。
通过上述保参数 hub 映射，可以研究本问题的对应下界。

论文表达时应把两层证据分开：
外部条件下界属于引用文献；hub 嵌入的语义、编码长度和参数关系属于本项目
需要列明的归约证明。只有把这些前提接好，才可以陈述
“本模型具有相应 ETH 条件障碍”。
它不说明现有带标签 DP 的每一个宽松常数和位宽指数都是最优的，
不排除在特殊树形、少数实际费用类型或其他结构参数上更快，
也不构成墙钟性能下界。
本文给出了语义嵌入，但没有在此次文档工作中增加归约编码或标准机器成本的 Lean 证明。

本项目已有的严格 hub 三划分归约及其编码包与这条新参数讨论也应区分。
原 \#23 的语义、gamma/EOF 解析和长度界已有范围明确的交付；
剩余机器构造与验证成本是独立任务。
把高重数装箱的参数下界加入文献定位，不会自动完成那些机器层义务，
也不能以此重新打开原已交付接口的验收范围。

\subsection{对称性、多重集与代表族的已有工作}

Kaibel 与 Pfetsch 的 orbitopes \cite{kaibel2008}
从列置换作用的规范代表研究整数规划对称消除；
Diaconis 与 Simper \cite{diaconis2021}
讨论的双陪集与列联表对应，为式 \eqref{eq:orbit-convolution} 提供代数背景。
这些来源说明：去掉相同车辆的编号与按类型记录匹配不是一个新的总方法。
本项目应当报告具体算法如何处理树支持、精确流量、有符号费用与见证恢复，
以及相较带标签版本实际减少多少状态和分配。

D\"oring 等人的 IPEC 2025 工作 \cite{doering2025}
已经在其容量路由树分解 DP 中按路线特征计数，
并通过多维 fingerprint bin packing 描述局部路线的组合。
所以“把局部路线视作物品再合并”也有直接先例。
其模型中的非负资源和可删减重复边引理，与这里固定重数和有符号费用不同。
比较时应指出哪条引理不适用、当前算法采用什么替代条件，
而不是用模型不同作为充分的新颖性论证。

Bodlaender、Cygan、Kratsch 与 Nederlof 的 rank-based 方法
\cite{bodlaender2015} 通过连接关系的代数结构保留小代表集合。
本项目可以借鉴“对所有未来补全保留最优性”的规格，
但不能直接借用其状态数结论。
第 \ref{sec:interface} 节给出的明确资源兼容矩阵具有很高的秩，
说明任何尝试都必须提出更合适的分解或更受限的补全语言，
并在该新关系上证明秩界。
这也是负面结构定理的正面用途：尽早淘汰没有适用前提的算法口号。

\section{三项研究计划与可验收产物}\label{sec:plan}

\subsection{优先级及与既有 issue 的关系}

建议保留已有 \#19、\#21、\#22、\#23 的范围记录，
把以下工作作为新研究包挂到主线，而不是把所有未做事项追加为旧包未完成的理由。
三个包可以并行开始，但 A 的模式语义与 B 的接口语义应优先稳定，
C 再据此固定真正充分的联合类型。
实施阶段编号以配套路线图的 A0--A4、B0--B4、C0--C4 为唯一编号；
下文按证明依赖叙述工作顺序，不另建一套与其竞争的任务编号。
具体任务分别记录在
\href{https://github.com/huam54925-cpu/uav-routing-certification/issues/24}{\#24（A）}、
\href{https://github.com/huam54925-cpu/uav-routing-certification/issues/25}{\#25（B）}、
\href{https://github.com/huam54925-cpu/uav-routing-certification/issues/26}{\#26（C）}。
地址级成本收口与冷启动验证继续作为独立质量任务推进；
它们对于可复现性有必要，但不替代这里的创新目标。

\begin{longtable}{p{0.08\textwidth}p{0.25\textwidth}p{0.34\textwidth}p{0.22\textwidth}}
\toprule
包 & 主要研究问题 & 最小可验收结果 & 应避免的过度承诺\\
\midrule
\endfirsthead
\toprule 包 & 主要研究问题 & 最小可验收结果 & 应避免的过度承诺\\ \midrule
\endhead
A &
一般二进制实例有没有短压缩证书？ &
$P_T$ 整数点等价；两种支持界及同一证书推论；与已有证书结构对接；
确定压缩编码、解析与位成本义务。 &
不能因短证书存在就宣称已经完成标准机器 NP 形式化。\\
B &
小树接口究竟必须保留哪些费用信息？ &
五顶点非负族；带标签与置换版兼容矩阵；
补全关系及代表集合下界；明确正向拼接契约。 &
不能把等价类数或矩阵秩说成任意算法、任意符号编码的空间下界。\\
C &
如何删掉车辆编号冗余而不丢可行匹配？ &
联合类型充分性；运输表可实现性；轨道 DP 完备性；
规范证书提升；类型计数及消融实验设计。 &
不能只排序后相加；不能默认保留旧 tie-breaking 或得到 FPT。\\
\bottomrule
\end{longtable}

\subsection{A 包：复用检查器，把创新集中在支持界}

仓库已有 \repo{Certificate} 与 \repo{checkCertificate}，
以及相应的 soundness/completeness 和模式重数语义，
现有检查并不要求先展开所有重复行。
因此 A 包不应为了新论文重新造一个同义的压缩证书类型。
应先确认原结构允许的零重数与空列表行为，
用规范化桥接把数学证书送入已有检查器。
尤其是原检查器可能仍检查零重数条目的模式；
删除一个无效但零重数条目可能改变任意原始输入的布尔结果。
“压缩保留有效证书语义”不等于“压缩保持所有原始字节输入的检查结果”。

先交付单行整数线性化与多面体凸性，
随后交付有限最小支持和平方势存在性所需引理，
把对数界、奇偶界与共同最小支持串起来；
然后选择规范二进制格式，复用原 checker 的语义，
补齐解析、数值位宽、循环次数和不展开重数的时间分析。
这些阶段的文档应分别列出具体定理、依赖和未完成接口。
测试重点不是重复证明有限计数恒等式，而是零边流、零路线、
有符号费用、不可约与未约分有理数、巨大重数及非法解析边界。

数学验收要求同时保留两个支持界，而不是只选一个较容易实现的上界。
但算法验收不要求立即找到最小支持证书；
存在一个有界支持证书和高效计算最小支持本来就是不同问题。
若实施时只能先证明存在性，应明确把构造算法留在下一阶段。

\subsection{B 包：先固定补全语言，再谈最小接口}

B 包首先应把“内部”“外部”“按标签拼接”“允许置换拼接”
定义为明确谓词，证明五顶点所有实例具有相同结构接口。
随后证明费用等价化为坐标不等式，利用固定总和得到单位兼容矩阵。
这条证明结构短、错误面小，适合作为优先形式化的非平凡结果。
同时记录所有费用非负、树大小固定、$m=2$ 即可产生大量等价类等性质，
使下界能够精准回应常见压缩建议。

接着把足够接口写成 gluing 定理：给定两侧局部证书和兼容匹配，
可以生成全局矩阵，并保持精确列和、归一化、支持和费用。
反向应说明任意全局矩阵在该割上能分解回被定义的两侧接口。
若只给正向构造，不足以支撑完整 DP；若只给反例，
又不足以指导应该存哪些字段。两个方向应被同一组定义连接。

本包验收不以“证明所有可能压缩都不行”为目标。
更合理的成功标准是：精确指出一类摘要或代表删除不可行，
同时给出下一步算法可以依赖的充分信息接口。
若后续发现一个符号化整数锥表示绕过显式代表表，
它与本包下界并不矛盾，反而说明下界范围界定正确。

\subsection{C 包：先证明商递推，再决定数据结构}

C 包的第一实现应优先可解释性。
保留旧带标签执行器作为语义参照，新增规范联合类型、
运输表枚举和匹配见证，不立即同时更换所有堆布局与排序算法。
先证明每次原合并都落入某个运输表结果，
每次运输表合并又能提升回原合并；对树递归归纳后得到根最优值相等。
再根据实际瓶颈选择稀疏表、共享 payload 或更快运输表生成方式。

实验应分开记录状态数量下降与单状态处理成本。
至少包含四类实例：大量重复费用与支持类型；
几乎每行费用都不同；负费用导致中间费用反复变化；
五顶点下界族或其可实现变体。
比较指标应包括原状态数、轨道数、实际类型数、运输表数、
规范化时间、证书提升时间和最终检查时间。
如果状态减少但匹配枚举过重，应如实报告；
如果仅在对称实例上有效，应给出与类型分布有关的参数说明。

新 API 的首版验收是最优值等价和原检查器接受，
另加稳定规范输出；旧完整证书的逐字段相等放在独立可选目标中。
最终层最大流或反序匹配可以作为受条件保护的特化，
其适用条件必须由程序结构或证明项提供，不能由调用者口头保证。
在得到正确的商递推后，再为它建立新的标量或地址成本模型，
避免对一个尚未证明完备的启发式过早投入精细计费。

\subsection{证明、工程与实验三条线如何互相支持}

建议每个结果都提供三项互不替代的信息：
数学命题及完整假设；形式化或实现状态；验证与实验的准确范围。
纸面定理可以先发表在研究记录里，但不能加上 Lean 已证的标签；
机器检查通过的函数可以先作为正确性基线，但不因此宣称最佳性能；
实验显著加速可以指导后续优化，却不能代替完备性证明。

冷启动构建、统一审计索引、版本固定和文档导航是必要收口任务。
它们应在投稿前形成可重现包，但不必把研究时间全部消耗在继续扩展底层执行模型。
对于成本模型，优先补会影响主要定理成立与否的缺口，
对只改变保守常数的精化可在研究结果明确后安排。
这能同时保持可信度与研究产出，而不是在“只做数学”和“只做工程”之间二选一。

\section{Lean 关键接口草案：未编译的设计片段}\label{sec:lean}

本节仅给短代码来说明证明应落在哪里，不构造伪完成的 Lean 文件。
片段没有在本轮运行工具链；真实命名、类型转换和导入需要本地核对。
没有使用 \repo{sorry} 或额外公理补洞。
已有 \repo{Certificate} 的结构字段应从仓库复用，下面不重复定义。

单行线性化的关键自然数事实可以独立写成如下引理。
它把“父边为正”转成一个无布尔变量的线性整数约束：
\begin{lstlisting}
import Mathlib

-- Design fragment only; not compiled in this document update.
theorem parent_guard_iff
    (x y cap : Nat) (hx : x <= cap) :
    x <= cap * y <-> (0 < x -> 0 < y) := by
  constructor
  · intro hxy hpos
    by_contra hnot
    have hy : y = 0 := Nat.eq_zero_of_not_pos hnot
    simp only [hy, Nat.mul_zero] at hxy
    omega
  · intro h
    by_cases hz : x = 0
    · simp [hz]
    · have hy : 1 <= y := h (Nat.pos_of_ne_zero hz)
      calc
        x <= cap := hx
        _ = cap * 1 := (Nat.mul_one cap).symm
        _ <= cap * y := Nat.mul_le_mul_left cap hy
\end{lstlisting}

证明奇偶界时，需要把“有限可行矩阵上存在最小平方势”与
“不同同奇偶行的中点仍合法”分成两个引理；
证明对数界时，应该先写生成向量子集和的注入性，
再调用有限集合基数界。这样不会把选择最小支持与选择最小平方势混为一个优化对象。
建议待实现接口名称与用途如下，名称只是设计建议：
\begin{center}
\begin{tabular}{p{0.43\textwidth}p{0.48\textwidth}}
\toprule
建议接口 & 数学契约\\ \midrule
\repo{row_polytope_integer_iff} & 整数点等价于已有合法行谓词。\\
\repo{exists_certificate_log_support} & 存在满足目标盒子对数界的有效证书。\\
\repo{exists_certificate_parity_support} & 存在支持不超过 $2^t$ 的有效证书。\\
\repo{exists_certificate_both_support} & 用全局最小支持同时取得两个界。\\
\repo{five_vertex_compatible_iff} & 带标签兼容恰好是向量相等。\\
\repo{five_vertex_orbit_compatible_iff} & 允许置换时兼容恰好是多重集相同。\\
\repo{coupling_realizable} & 每张满足边际的运输表可提升为具体排列。\\
\repo{orbit_solver_optimum_eq} & 新 API 与旧 API 的最优值相等。\\
\bottomrule
\end{tabular}
\end{center}

运输表的数据契约可以使用一个很小的结构表示；以下也只是未编译片段：
\begin{lstlisting}
import Mathlib
open scoped BigOperators

structure RowCoupling
    (A B : Type*) [Fintype A] [Fintype B]
    (mu : A -> Nat) (nu : B -> Nat)
    (compatible : A -> B -> Prop) where
  count : A -> B -> Nat
  row_sum :
    forall a, Finset.univ.sum (fun b => count a b) = mu a
  col_sum :
    forall b, Finset.univ.sum (fun a => count a b) = nu b
  forbidden :
    forall a b, Not (compatible a b) -> count a b = 0
\end{lstlisting}

上面使用 \repo{Finset.univ.sum} 明确有限索引域；
接入仓库时仍需核对实际类型类、命名空间与现有联合类型定义，
不能把本稿未编译的展示代码当作已验证成品。
关键证明输出不只是一个布尔真值，而是一个
$\operatorname{Equiv.Perm}(\operatorname{Fin} m)$
及其联合计数等式。这个排列随后必须同时作用于结构字段、费用字段与矩阵 payload。

对于二进制检查器，应首先把一个成功解析的压缩列表连接到现有
\repo{checkCertificate}，再证明被接受时输出语义正确。
拒绝所有非规范编码、对巨大自然数的早期守卫和中间位宽，
属于编码执行层，不适合混进最初的有限组合引理。
分层的目的不是减少证明要求，而是使每个公开定理的输入、输出及成本边界明确。

\section{如何评价现有工作并组织投稿主线}\label{sec:evaluation}

“工作量很大”与“数学原创突破”是两个不同维度。
现有仓库的强项在于语义链长且细节多：
精确访问、负费用、完整证书顺序、编码、拒绝分支和计费范围都受到认真处理。
这使它可以支持有分量的形式化研究。
但仅凭模块数、公理审计或几组回归，不能预先保证任何会议或期刊接收；
论文仍要说明技术难点、相对于文献的新增内容和可复现证据。

外部评价中关于成本模型、清洁构建和文档组织的提醒值得保留。
不过“因此只能继续补工程”并不是唯一结论。
本稿三项研究包就在现有稳定语义上提出新的结构问题，
并允许用较短的数学证明先验证其价值。
其中 A 解决压缩证书长度，B 给出信息保留的必要性，
C 对准算法重复枚举；三者可以形成一条互相解释的论文线。

也不应因为当前模型局限于树就否定应用来源。
合理表述是：UAV 场景促成了精确重复访问和认证需求，
而当前论文研究其可独立陈述的树支持子问题。
若主要定理只涉及这个子问题，题目、摘要与实验就应围绕它，
而不要暗示已经解决一般 UAV 调度。
未来回到一般图、时间窗或物理轨迹时，
应首先寻找保持当前证书语义的保真嵌入或松弛认证，
而不是同时把所有应用特征塞进现有状态。

建议主论文的问题陈述先回答三件事：
为什么精确边流与逐次 signed 费用使常见简化失效；
为什么仍能得到短模式证书及可认证求解；
为什么某些状态压缩不可能而置换商压缩仍然正确。
带地址执行模型可作为充分支撑可信实现的贡献，
但不必让每一个历史计费包都成为正文同等篇幅的主角。
具体常数、审计清单、版本迁移与完整实现索引适合放在附录或配套材料。

投稿前至少要把每条主张放到下面一种明确类别：
\begin{enumerate}
\item \textbf{当前 Lean 已证：}可链接到固定版本的定理及审计入口，
      同时列出未被消除的树数据、非空性或成本模型假设。
\item \textbf{本文纸面已论证：}如线性化、双支持界、五顶点接口族、
      运输表完整性，给出可复核证明但不宣称已经通过内核。
\item \textbf{文献推论：}明确外部定理版本、参数代入及编码条件，
      不将通用算法的原创性归于本项目。
\item \textbf{待实施或待验证：}商 DP、压缩二进制检查时间、
      新实验与统一冷构建，列出验收条件和失败后的回退方案。
\end{enumerate}
这样的划分不是降低研究价值，而是使真正新增的结果更突出。
例如一个完整的五顶点单位矩阵族，比含糊声称“状态不能再压缩”更有内容；
一个保持检查器语义的运输表商递推，比声称“做了对称消除所以快很多”更可审查。
清晰的限制也能避免项目被一个过大的“端到端”标签绑住。

\section{本稿的证据清单与下一次更新方式}

本稿没有改变任何 Lean 模块、数据集、工具链或依赖锁定文件。
所有现有正式成果的概述以已核对远端基线和其状态文档为准；
本稿新增的定理是纸面论证，不是新的构建记录。
LaTeX 源码自包含参考文献，不依赖外部图片或 bib 文件，
但没有在本轮编译，不能据此报告页数、无溢出或 PDF 检查通过。
本地可使用 XeLaTeX；若调整字体或版面，应单独记录排版检查，不与数学验证混同。

后续每次更新应包含：新增或修改的命题；
其依赖的仓库基线；证明从纸面推进到 Lean 的具体状态；
是否改变算法规格；是否产生新的实验与其可复现入口。
如果一个计划因反例而失败，应保留反例和修订后的范围，
而不是继续使用原来的过强标题。
如果某个数学推论已被近期文献覆盖，应及时把新颖性主张转到模型化、证书接口、
形式化或实证改进上。

本次最合理的下一步，是先实现 A 的两个短支持证明与 B 的五顶点族，
用它们锁定证书规模和接口语义；随后开发 C 的规范匹配与轨道递推。
旧的完整证书求解器继续作为语义基准，
其计费与历史验证结果保持原范围。这样项目可以在已有工作上增加可识别的结构成果，
同时保留能够审查、复现和逐步扩展的证明基础。

\begin{thebibliography}{99}

\bibitem{repository}
UAV Routing Certification 项目.
\emph{研究分支及固定版本的 README、预算查询、堆去重与路线输出说明}.
基线 \repo{f3ea0823be69ffa7152dec903069c2ba79ad20a4}, 2026-11-18.
\url{https://github.com/huam54925-cpu/uav-routing-certification}.
配套文献审查见
\repo{docs/TREE_BUDGET_LITERATURE_REVIEW_20261118.md}，
实施路线见
\repo{docs/TREE_BUDGET_INNOVATION_ROADMAP_20261118.md}。

\bibitem{eisenbrand2006}
F. Eisenbrand and G. Shmonin.
Carath\'eodory bounds for integer cones.
\emph{Operations Research Letters}, 34(5):564--568, 2006.
\url{https://doi.org/10.1016/j.orl.2005.09.008}.

\bibitem{goemans2014}
M. X. Goemans and T. Rothvo\ss.
Polynomiality for Bin Packing with a Constant Number of Item Types.
In \emph{Proceedings of SODA 2014}, 830--839, 2014.
\url{https://arxiv.org/abs/1307.5108}.
本文采用该版本 Theorem 2.2 的整数锥表述。

\bibitem{goemans2020}
M. X. Goemans and T. Rothvo\ss.
Polynomiality for Bin Packing with a Constant Number of Item Types.
\emph{Journal of the ACM}, 67(6), Article 38, 2020.
\url{https://doi.org/10.1145/3421750}.

\bibitem{jansen2026convex}
K. Jansen, F. Ohnesorge, and B. Schumacher.
\emph{A Constructive FPT Algorithm for Integer Cone Optimization with Convex Costs}.
arXiv:2609.31328, version 5, 2026-11-03.
\url{https://arxiv.org/abs/2609.31328v5}.
本稿以已核对的第 5 版 PDF 的题名及 Theorem 1 为准；
检索索引中可能仍显示早期题名。

\bibitem{jansen2026lower}
K. Jansen, F. Ohnesorge, and L. Pirotton.
A Tight Double-Exponential Lower Bound for High-Multiplicity Bin Packing.
In \emph{ICALP 2026}, LIPIcs 374, 116:1--116:20, 2026.
\url{https://doi.org/10.4230/LIPIcs.ICALP.2026.116}.

\bibitem{kaibel2008}
V. Kaibel and M. E. Pfetsch.
Packing and Partitioning Orbitopes.
\emph{Mathematical Programming}, 114:1--36, 2008.
\url{https://arxiv.org/abs/math/0603678}.

\bibitem{diaconis2021}
P. Diaconis and M. Simper.
\emph{Statistical Enumeration of Groups by Double Cosets}.
arXiv:2102.04576, 2021.
\url{https://arxiv.org/abs/2102.04576}.

\bibitem{doering2025}
M. D\"oring, J. Fehse, T. Friedrich, P. Marten, N. Mohrin,
K. Simonov, F. Soheil, J. Timm, and S. Verma.
Parameterized Complexity of Vehicle Routing.
In \emph{IPEC 2025}, LIPIcs 358, Article 10, 2025.
\url{https://doi.org/10.4230/LIPIcs.IPEC.2025.10}.
本文主要比较其第 4.4 节的路线计数接口与合并方法。

\bibitem{bodlaender2015}
H. L. Bodlaender, M. Cygan, S. Kratsch, and J. Nederlof.
Deterministic single exponential time algorithms for connectivity problems
parameterized by treewidth.
\emph{Information and Computation}, 243:86--111, 2015.
预印本题名为
\emph{Solving weighted and counting variants of connectivity problems
parameterized by treewidth deterministically in single exponential time}.
\url{https://arxiv.org/abs/1211.1505}.

\end{thebibliography}
\end{document}
